% BÀI TOÁN HÌNH HỌC - BAI_HINH_36 (TRÍCH ĐỀ 36)
\begin{ex}
% Câu 13
	Cho tam giác ABC nhọn (AB < AC) nội tiếp đường tròn (O;R). Các đường cao AK, BE, CD cắt nhau tại H.
	\begin{enumerate}[a)]
		\item Chứng minh bốn điểm B, D, E, C cùng thuộc một đường tròn.
		\item Gọi N là giao điểm của DC và KE, P là hình chiếu vuông góc của C lên đường kính AM. Khi A, B cố định, C thay đổi nhưng vẫn thỏa mãn bài toán, chứng minh NH.CD = DH.NC và xác định vị trí của C để chu vi tam giác OCP lớn nhất.
	\end{enumerate}
\end{ex}
